\appendix
\chapter{Trazabilidad de par\'ametros de modelado}
\label{anexo:trazabilidad}

La \cref{tab:anexo_trazabilidad} recopila la totalidad de constantes f\'isicas y de control parametrizadas en \texttt{backend/config.py}, se\~nalando su valor num\'erico, origen metodol\'ogico y documento de soporte seg\'un el documento 06 (\S 3).

\begin{longtable}{@{}L{4.3cm}L{2.5cm}L{2.3cm}L{6.2cm}@{}}
\caption{Matriz exhaustiva de trazabilidad de par\'ametros del modelo Axioma-H.}
\label{tab:anexo_trazabilidad} \\
\toprule
\textbf{Constante} & \textbf{Valor nominal} & \textbf{Categor\'ia} & \textbf{Fuente / Justificaci\'on} \\
\midrule
\endfirsthead
\toprule
\textbf{Constante} & \textbf{Valor nominal} & \textbf{Categor\'ia} & \textbf{Fuente / Justificaci\'on} \\
\midrule
\endhead
\bottomrule
\endfoot
\texttt{UMBRAL\_TEMPERATURA\_CORTE\_C} & $45.0\;\si{\celsius}$ & PRD / D1 & L\'imite t\'ermico absoluto (REGLA 01). \\
\texttt{TEMPERATURA\_OBJETIVO\_PULSO\_C} & $44.0\;\si{\celsius}$ & Matriz / D1 & Consigna t\'ermica de regeneraci\'on (REGLA 03). \\
\texttt{UMBRAL\_HUMEDAD\_ECOLOGICA\_PCT} & $30.0\,\%$ & PRD & L\'imite para suspensi\'on por sequ\'ia (REGLA 02). \\
\texttt{UMBRAL\_RESERVA\_CRITICA\_PCT} & $30.0\,\%$ & PRD / D2 & L\'imite de bater\'ia para inhibir pulso (REGLA 04). \\
\texttt{UMBRAL\_IMPEDANCIA\_OHMS} & $1500.0\;\Omega$ & PRD & Detecci\'on de cristalizaci\'on salina. \\
\texttt{IMPEDANCIA\_NOMINAL\_OHMS} & $1000.0\;\Omega$ & PRD & Resistencia superficial de membrana limpia. \\
\texttt{DURACION\_PULSO\_SEG} & $90.0\;\text{s}$ & PRD & Duraci\'on del calentamiento resistivo. \\
\texttt{FRACCION\_MAXIMA\_EXTRACCION} & $0.33$ ($33.0\,\%$) & Matriz / D3 & Tope de extracci\'on ecol\'ogica (Conjetura 3). \\
\texttt{CAUDAL\_DURANTE\_PULSO\_PCT} & $80.0\,\%$ & PRD & Caudal permitido durante pulso t\'ermico. \\
\texttt{CAUDAL\_RESERVA\_CRITICA\_PCT} & $50.0\,\%$ & PRD & Caudal restringido sin pulso (REGLA 04). \\
\texttt{CAUDAL\_NOMINAL\_PCT} & $100.0\,\%$ & PRD & Operaci\'on nominal continua (REGLA 05). \\
\texttt{RECUPERACION\_MINIMA\_PCT} & $90.0\,\%$ & Matriz & Umbral de \'exito de Conjetura 1. \\
\texttt{AHORRO\_ENERGIA\_MINIMO\_PCT} & $35.0\,\%$ & Matriz & Umbral de \'exito energ\'etico de Conjetura 1. \\
\texttt{TIEMPO\_REGENERACION\_MAX\_SEG} & $120.0\;\text{s}$ & Matriz & Umbral temporal cin\'etico de Conjetura 2. \\
\texttt{FACTOR\_IMPEDANCIA\_REGENERADA} & $1.10$ ($1100\;\Omega$) & Criterio equipo & Criterio de membrana limpia tras pulso. \\
\texttt{PERDIDA\_TERMICA\_MAXIMA\_PCT} & $0.0\,\%$ & Matriz & Tolerancia cero de degradaci\'on t\'ermica (C2). \\
\texttt{VIOLACIONES\_ECOLOGICAS\_MAXIMAS} & $0$ & Matriz & Tolerancia cero de extracci\'on con $\text{HR}<30\,\%$ (C3). \\
\texttt{RENDIMIENTO\_MINIMO\_L\_M2\_NOCHE} & $12.0\;\text{L/m}^2\text{/noche}$ & Matriz & Umbral de rendimiento de Conjetura 3. \\
\texttt{FLUJO\_NIEBLA\_MAXIMO\_L\_M2\_H} & $6.0\;\text{L/m}^2\text{/h}$ & Supuesto (D5) & Flujo saturado l\'inea base; optimista vs. literatura. \\
\texttt{FLUJO\_NIEBLA\_CALIBRADO\_MIN\_L\_M2\_H} & $0.49\;\text{L/m}^2\text{/h}$ & Calibraci\'on & Inversi\'on num\'erica a $3\;\text{L/m}^2\text{/d\'ia}$ \parencite{klemm2012fog}. \\
\texttt{FLUJO\_NIEBLA\_CALIBRADO\_MAX\_L\_M2\_H} & $1.83\;\text{L/m}^2\text{/h}$ & Calibraci\'on & Inversi\'on num\'erica a $10\;\text{L/m}^2\text{/d\'ia}$ \parencite{klemm2012fog}. \\
\texttt{EFICACIA\_MAXIMA\_PULSO} & $0.90$ & Supuesto & Fracci\'on m\'axima de ensuciamiento removido. \\
\texttt{CARGA\_SALINA\_REFERENCIA\_OHMS} & $2500.0\;\Omega$ & Supuesto & Escala de saturaci\'on de remoci\'on por pulso. \\
\texttt{CONSTANTE\_TIEMPO\_DISOLUCION\_SEG}& $25.0\;\text{s}$ & Supuesto & Constante cin\'etica $\tau_0$ a saturaci\'on t\'ermica. \\
\texttt{CAPACIDAD\_TERMICA\_MATRIZ\_J\_K} & $800.0\;\text{J/K}$ & Supuesto & Inercia t\'ermica de la superficie colectora. \\
\texttt{CONDUCTANCIA\_PERDIDAS\_W\_K} & $30.0\;\text{W/K}$ & Supuesto & Coeficiente de convecci\'on y disipaci\'on t\'ermica. \\
\texttt{CONSUMO\_BASE\_NEURO\_W} & $0.6\;\text{W}$ & Supuesto & Potencia de electr\'onica y sensores de reglas. \\
\texttt{CONSUMO\_BASE\_TEMPORIZADOR\_W} & $0.2\;\text{W}$ & Supuesto & Potencia el\'ectrica base de circuito temporizador. \\
\texttt{INTERVALO\_PULSO\_TEMPORIZADOR\_MIN}& $120.0\;\text{min}$ & Supuesto & Per\'iodo invariable de activaci\'on del temporizador. \\
\texttt{FATIGA\_POR\_PULSO\_PCT} & $\SI{0.02}{\percent}$ & Supuesto & P\'erdida mec\'anica por ciclado t\'ermico suave. \\
\texttt{DEGRADACION\_TERMICA\_PCT\_POR\_C\_MIN} & $\SI{0.05}{\percent\per(\celsius\cdot\minute)}$ & Supuesto & Tasa de da\~no por exceso t\'ermico ($>45\;\si{\celsius}$). \\
\texttt{DEGRADACION\_CRISTALIZACION\_PCT\_H}& $\SI{0.10}{\percent\per\hour}$ & Supuesto & Da\~no por permanencia prolongada de salitre. \\
\texttt{CONSUMO\_SUPERVIVENCIA\_L\_PERSONA\_DIA} & $15.0\;\text{L/p/d\'ia}$ & Referencia & Est\'andar humanitario Sphere \parencite{sphere2018handbook}. \\
\texttt{CONSUMO\_ACCESO\_BASICO\_L\_PERSONA\_DIA} & $20.0\;\text{L/p/d\'ia}$ & Referencia & Est\'andar acceso b\'asico OMS \parencite{howard2020domestic}. \\
\texttt{RENDIMIENTO\_LITERATURA\_MIN\_L\_M2\_DIA} & $3.0\;\text{L/m}^2\text{/d\'ia}$ & Referencia & M\'inimo operacional t\'ipico de \textcite{klemm2012fog}. \\
\texttt{RENDIMIENTO\_LITERATURA\_MAX\_L\_M2\_DIA} & $10.0\;\text{L/m}^2\text{/d\'ia}$ & Referencia & M\'aximo operacional t\'ipico de \textcite{klemm2012fog}. \\
\texttt{ENTREGA\_CARROTANQUE\_L\_COMUNIDAD\_MES} & $1000.0\;\text{L/com/mes}$ & Referencia & Asignaci\'on mensual asistencial ICBF \parencite{icbf2025agua}. \\
\texttt{INCRUSTACION\_TRANSICION\_MIN\_OHM\_DIA} & $1171.06\;\Omega\text{/d\'ia}$ & Modelo & L\'imite inferior de la banda de cruce (semilla 42). \\
\texttt{INCRUSTACION\_TRANSICION\_MAX\_OHM\_DIA} & $2365.34\;\Omega\text{/d\'ia}$ & Modelo & L\'imite superior de la banda de cruce (semilla 42). \\
\end{longtable}

\chapter{Diccionario de datos abiertos de investigaci\'on}
\label{anexo:datos_abiertos}

El repositorio publica los datos de referencia en formato abierto bajo el est\'andar Frictionless Data Package (\texttt{docs/datos/datapackage.json}). El conjunto comprende cinco archivos tabulares:
\begin{enumerate}[leftmargin=*]
  \item \textbf{\texttt{barrido\_salinidad.csv} (6 registros):} describe el comportamiento promediado en cinco semillas estoc\'asticas para cargas salinas de 20 a $240\;\Omega\text{ por L/m}^2$. Campos principales: recuperaci\'on media y desviaci\'on de ambos controles (\texttt{recuperacion\_neuro\_media}, \texttt{recuperacion\_temporizador\_media}), degradaci\'on de integridad acumulada (\texttt{degradacion\_neuro\_media}, \texttt{degradacion\_temporizador\_media}), porcentaje de reducci\'on de desgaste (\texttt{reduccion\_desgaste\_pct}), tasa de incrustaci\'on diaria (\texttt{incrustacion\_diaria\_ohm\_media}) y ventaja neta en puntos porcentuales (\texttt{ventaja\_recuperacion\_pts}).
  \item \textbf{\texttt{robustez\_respuesta.csv} (54 registros):} matriz factorial completa combinando flujos de niebla ($6.0$, $1.83$, $0.49\;\text{L/m}^2\text{/h}$), eficacias de pulso ($0.80$, $0.90$, $1.00$) y las 6 tasas de salinizaci\'on. Incluye el valor del cruce interpolado (\texttt{cruce\_incrustacion\_diaria\_ohm}) y los niveles de la cuadr\'icula entre los que se ubica (\texttt{cruce\_incrustacion\_inferior\_ohm}, \texttt{cruce\_incrustacion\_superior\_ohm}).
  \item \textbf{\texttt{robustez\_desgaste.csv} (18 registros):} evaluaci\'on de sensibilidad mec\'anica variando fatiga por pulso ($0.01\,\%$, $0.02\,\%$, $0.04\,\%$) y degradaci\'on t\'ermica ($0.025$, $0.05$, $0.10\,\%/(\si{\celsius}\cdot\text{min})$). Registra la reducci\'on m\'inima y m\'axima de desgaste observada (\texttt{reduccion\_desgaste\_min\_pct}, \texttt{reduccion\_desgaste\_max\_pct}) y la fracci\'on de escenarios favorables (\texttt{puntos\_con\_menor\_desgaste}/\texttt{puntos\_totales}).
  \item \textbf{\texttt{impacto\_entornos.csv} (18 registros):} matriz de dimensionamiento comunitario que cruza 3 reg\'imenes de niebla por 6 cargas salinas. Registra el rendimiento diario (\texttt{rendimiento\_neuro\_l\_m2\_dia}), clasificaci\'on frente al cruce (\texttt{zona\_incrustacion}), colectores requeridos para 50 personas a $20\;\text{L/d\'ia}$ (\texttt{nodos\_requeridos}) y nodos necesarios para igualar los $1000\;\text{L/mes}$ del carrotanque (\texttt{nodos\_equivalentes\_entrega}).
  \item \textbf{\texttt{resultados\_escenarios.csv} (10 registros):} resultados detallados de la simulaci\'on en los diez escenarios de operaci\'on y perturbaci\'on clim\'atica (sequ\'ia, tormenta de arena, calor extremo y combinaciones) con humedad inicial de $85\,\%$ y $45\,\%$.
\end{enumerate}

\chapter{Mapa de evidencias frente a criterios de evaluaci\'on}
\label{anexo:mapa_evidencias}

El \cref{tab:mapa_evidencias} vincula de manera objetiva los diez criterios de evaluaci\'on del curso de investigaci\'on con las secciones de este informe y los m\'odulos de c\'odigo fuente del repositorio, facilitando la verificaci\'on acad\'emica independiente.

\begin{table}[htbp]
\centering
\small
\caption{Mapeo objetivo de evidencias t\'ecnicas frente a criterios de evaluaci\'on acad\'emica.}
\label{tab:mapa_evidencias}
\begin{tabularx}{\textwidth}{@{}L{2.8cm}L{4.2cm}Y@{}}
\toprule
\textbf{Criterio} & \textbf{Secci\'on del informe} & \textbf{Archivo o artefacto en el repositorio} \\
\midrule
\textbf{Cobertura} & Cap. \ref{cap:planteamiento}, \ref{cap:metodologia}, \ref{cap:diseno_artefacto}, \ref{cap:resultados} & \texttt{docs/01}, \texttt{docs/02}, \texttt{backend/main.py}, \texttt{docs/api/openapi.json} \\
\textbf{Coherencia} & Cap. \ref{cap:metodologia}, \ref{cap:resultados}, \ref{cap:discusion_calidad} & \texttt{docs/04} (\S 7, \S 8), \texttt{tests/test\_contenido.py}, \texttt{frontend/contenido.py} \\
\textbf{Profundidad} & Cap. \ref{cap:metodologia}, \ref{cap:resultados} & \texttt{backend/analisis.py}, \texttt{docs/datos/*.csv}, \texttt{docs/04} (\S 8.6 a \S 8.11) \\
\textbf{S\'intesis} & Resumen, Lectura breve, Cap. \ref{cap:resultados} & \texttt{README.md}, \texttt{frontend/contenido.py} (\texttt{VEREDICTOS}, \texttt{LECTURA\_BREVE}) \\
\textbf{Originalidad} & Cap. \ref{cap:marco_teorico}, \ref{cap:diseno_artefacto}, \ref{cap:conclusiones} & Arquitectura de reglas (\texttt{backend/motor\_reglas.py}), decisi\'on D5 \\
\textbf{Fluidez} & Estructura global del informe & Formato \LaTeX{} con hiperv\'inculos internos, tablas Booktabs y gr\'aficas Pgfplots \\
\textbf{Relevancia} & Cap. \ref{cap:introduccion}, \ref{cap:planteamiento}, \ref{cap:conclusiones} & Contexto territorial de La Guajira, Sentencia T-302, est\'andares OMS/Sphere \\
\textbf{Claridad} & Lectura breve, Glosario, Cap. \ref{cap:diseno_artefacto} & Selector de niveles de lectura en Streamlit (\texttt{frontend/app.py}) \\
\textbf{Confiabilidad} & Cap. \ref{cap:discusion_calidad}, Anexo \ref{anexo:trazabilidad} & 828 pruebas automatizadas, 100\,\% cobertura, SonarQube Cloud, Render CI \\
\textbf{Universalidad} & Cap. \ref{cap:conclusiones}, Anexo \ref{anexo:guias_lectura} & Desacoplamiento REST, formato abierto Frictionless Data, gu\'ia interdisciplinaria \\
\bottomrule
\end{tabularx}
\end{table}

\chapter{Gu\'ias de lectura del artefacto por p\'ublico y disciplina}
\label{anexo:guias_lectura}

Para facilitar la navegaci\'on del software y de la documentaci\'on t\'ecnica seg\'un el perfil del usuario, el proyecto define trayectorias de lectura diferenciadas en el tablero y en los documentos.

\section{Gu\'ia por tipo de p\'ublico}
El \cref{tab:publicos} orienta la lectura en funci\'on del rol del evaluador o usuario.

\begin{table}[htbp]
\centering
\small
\caption{Recomendaci\'on de itinerarios de lectura seg\'un el perfil del usuario.}
\label{tab:publicos}
\begin{tabularx}{\textwidth}{@{}L{3.5cm}L{3.2cm}Y@{}}
\toprule
\textbf{P\'ublico objetivo} & \textbf{Nivel sugerido} & \textbf{Secciones recomendadas} \\
\midrule
\textbf{Comunidad local y l\'ideres} & Lectura breve & Secci\'on preliminar de Lectura Breve; Vistas 1 y 4 (parte d: impacto comunitario) del tablero interactivo. \\
\textbf{Evaluadores de pol\'itica y gestores} & Lectura breve & Resumen ejecutivo, Veredictos (\cref{tab:veredictos}); Vistas 1 y 4 del tablero; secci\'on Respuesta en \texttt{README.md}. \\
\textbf{Investigadores y desarrolladores} & Lectura completa & Informe completo (\LaTeX{}); Vistas 4 y 5 del tablero; documento 04 (\S 7 y \S 8) y datos abiertos en \texttt{docs/datos/}. \\
\bottomrule
\end{tabularx}
\end{table}

\section{Gu\'ia de lectura interdisciplinaria}
El \cref{tab:disciplinas} detalla c\'omo abordar el artefacto desde cuatro campos de conocimiento distintos, articulando las vistas del software con preguntas cient\'ificas de cada \'area.

\begin{table}[htbp]
\centering
\small
\caption{Gu\'ia de exploraci\'on interdisciplinaria del artefacto Axioma-H.}
\label{tab:disciplinas}
\begin{tabularx}{\textwidth}{@{}L{3.2cm}L{4.2cm}Y@{}}
\toprule
\textbf{Disciplina} & \textbf{Vistas y par\'ametros clave} & \textbf{Pregunta cient\'ifica que permite explorar} \\
\midrule
\textbf{Hidrolog\'ia e Ingenier\'ia ambiental} & Vista 4 (partes b y d); Vista 5 (perfil de sitio y calibraci\'on). Campos: \texttt{rendimiento\_neuro\_media}, \texttt{incrustacion\_diaria} y \texttt{area\_colector\_m2}. & ?`El r\'egimen de niebla medido en un sitio espec\'ifico, calibrado con el modelo, genera con el aerosol marino local una incrustaci\'on superior al umbral de cruce de H1-Ciclo 3? \\
\textbf{Salud p\'ublica y Epidemiolog\'ia} & Vista 3 (impacto comunitario); Vista 4 (parte d: dimensionamiento). Campos: \texttt{personas\_acceso\_basico} y \texttt{nodos\_requeridos}. & ?`Cu\'antos captadores pasivos se requieren para abastecer de forma segura a una rancher\'ia de 50 personas bajo el est\'andar b\'asico OMS ($20\;\text{L/d\'ia}$)? \\
\textbf{Ecolog\'ia y Biolog\'ia de la conservaci\'on} & Vista 2 (pausa ecol\'ogica); Vista 3 (violaciones ecol\'ogicas y l\'imite del $33\,\%$). Campos: \texttt{violaciones\_ecologicas}. & ?`Limitar la extracci\'on al $33\,\%$ y pausar con $\text{HR} < 30\,\%$ previene el desecamiento del microclima en ecosistemas xerof\'iticos litorales? \\
\textbf{Ciencias sociales y Antropolog\'ia} & Vista 1 (perspectivas \'eticas y supuestos); Vista 4 (costo ecol\'ogico en personas abastecidas). & ?`C\'omo conciliar el principio precautorio del tope ecol\'ogico con las urgencias comunitarias, y qu\'e modelos de gobernanza comunitaria aseguran la autonom\'ia de la gesti\'on del agua? \\
\bottomrule
\end{tabularx}
\end{table}

